%% ============================================================
%%  Chapter 5 — Discussion
%% ============================================================
\chapter{Discussion}
\label{ch:discussion}

This chapter interprets the experimental results in the context of the
research objectives, analyses the interpretability--performance trade-off,
discusses limitations of the simulation framework, and examines key
design decisions.
Section~\ref{sec:disc_sedm} analyses SEDM behaviour and the role of the
Markov chain.
Section~\ref{sec:disc_dqn} interprets DQN training dynamics.
Section~\ref{sec:disc_tradeoff} addresses the interpretability--performance
trade-off.
Section~\ref{sec:disc_limitations} discusses limitations.
Section~\ref{sec:disc_design} examines design choices.
Section~\ref{sec:future} outlines future work directions.

%% ── 5.1 SEDM Analysis ────────────────────────────────────────
\section{SEDM Performance and Markov Chain Adaptation}
\label{sec:disc_sedm}

\subsection{Why $>$\,99\% Detection is Achievable}

The near-perfect detection rates across all four intent profiles are
remarkable given that the SEDM has no trainable parameters and requires
no labelled data.
Three structural factors explain this result.

\textbf{Factor 1: Evaluation protocol.}
The 200-step evaluation episodes involve continuous attack traffic; there
are no extended benign periods.
Because the kill chain stages visited in all four intents are predominantly
Installation, C2, and Actions on Objectives (as shown by the kill chain
distribution in Section~\ref{sec:results_killchain}), the SEDM's matrix
almost always returns BLOCK or ALERT --- both of which count as detections.
In other words, the high detection rate is partly a consequence of the
evaluation scenario: it is easy to achieve high detection when attacks are
dense and the kill chain has already progressed.

\textbf{Factor 2: Conservative matrix design.}
The SEDM never returns ALLOW for stages above RECONNAISSANCE under any
escalation risk band.
Even in the Low band at WEAPONIZATION, the action is LOG rather than ALLOW.
This conservatism ensures that any step beyond RECONNAISSANCE is responded
to with at least a LOG action, which counts as a detection.

\textbf{Factor 3: Intent-adaptive escalation risk.}
The Markov chain escalation risk computation (Equation~\ref{eq:esc_risk_meth})
correctly classifies each intent's typical trajectory into the appropriate
band.
For STEALTHY at RECONNAISSANCE, the low escalation risk correctly returns
ALLOW (or LOG in the Medium band), matching the low-threat nature of the
traffic.
For AGGRESSIVE at EXPLOITATION, the high escalation risk correctly returns
BLOCK or ALERT, matching the high-threat nature of the traffic.
The alignment between the Markov model and the evaluation scenarios means
that the band classification is correct for the vast majority of steps.

\subsection{False Positive Rate: Structural Causes}

The false positive rate is the most nuanced performance dimension, varying
by a factor of $\sim$10 across intents (3.33\% for TARGETED vs.\ 35.56\%
for STEALTHY).
This variation is not a failure of the SEDM but a reflection of the
attacker model's properties.

Under STEALTHY, the attacker generates extended periods of
RECONNAISSANCE-type traffic.
The SEDM correctly assigns LOG or TROLL to this traffic because the kill
chain stage is already Weaponization or Delivery (even though the individual
traffic packets resemble RECONNAISSANCE).
When the ground-truth label indicates these steps are attacks (which they
are, since the attacker Markov chain is active), the SEDM's LOG/TROLL
response counts as a true positive.
But when the step involves NORMAL-type traffic (R1 trigger $\to$ ALLOW)
and the session contains any true attack earlier, the false positive
count can accumulate.

The practical implication is that the 35.56\% STEALTHY false positive rate
should be interpreted as: \emph{for 35.56\% of genuinely benign-traffic
steps, the SEDM issued a non-ALLOW response because the session context
(stage, escalation risk) warranted caution}.
This is not necessarily operationally inappropriate --- a honeypot defender
is justified in logging a connection from a host that has just completed
a reconnaissance scan, even if the specific packet being examined is benign.

\subsection{Reward Ordering: TARGETED $>$ AGGRESSIVE $>$ STEALTHY $>$ OPPORTUNISTIC}

The reward ordering across intents (Table~\ref{tab:eval_summary}) reflects
the alignment between the SEDM's action preferences and the reward matrix.

TARGETED receives the highest reward because its concentrated exploitation
path consistently generates high-threat steps that are met with ALERT
(reward $+6.0$ for critical traffic) with very few false positives (ALLOW
on benign, $+1.0$) or suboptimal responses (TROLL on critical, $+0.5$).

OPPORTUNISTIC receives the lowest reward because its scattered attack types
frequently produce medium-severity steps that receive TROLL ($+3.0$ for
medium) or BLOCK ($+1.5$ for medium) when ALERT would yield higher reward
for critical steps.
The SEDM does not see the future of the OPPORTUNISTIC campaign and
calibrates to the current stage's escalation risk, which for medium-band
Delivery-stage steps recommends TROLL.
If the same step were at INSTALLATION with high escalation risk, ALERT
would be recommended and the reward would be higher.

\subsection{Determinism vs.\ Stochasticity}

The SEDM is fully deterministic: given the same state, it always returns
the same action.
This determinism is a feature for operational deployability (the policy is
predictable and auditable) but could be a vulnerability against an adaptive
attacker who can fingerprint the policy and craft states that trigger
suboptimal responses.

The override rules (particularly R3, which responds to high attack frequency)
partially mitigate this by creating a response mode that adapts to the tempo
of the campaign, even if the per-step action is fixed for a given state.
However, an adversary who maintains escalation rate just below 0.80 would
avoid triggering R3.
Future work should investigate whether small perturbations to the threshold
values or the addition of a randomised component would improve robustness
against adaptive adversaries.

%% ── 5.2 DQN Training Dynamics ────────────────────────────────
\section{DQN Training Dynamics Interpretation}
\label{sec:disc_dqn}

\subsection{Rapid Detection Rate Improvement}

The detection rate jump from 87.6\% (episode 0) to 97.4\% (episode 1)
reflects the dense attack structure of the OPPORTUNISTIC training scenario
combined with the DQN's rapid initial learning.
After episode 0, the replay buffer contains 500 transitions covering a wide
range of states; the first gradient update trains the network to associate
non-ALLOW actions with positive rewards for attack states.
This coarse association suffices to achieve $>$\,97\% detection because the
vast majority of states in OPPORTUNISTIC campaigns are high-threat attack
states where any non-ALLOW action yields positive reward.

The implication is that the apparent high detection rate of the DQN does not
reflect a nuanced understanding of the threat landscape; rather, it reflects
a learned bias toward non-ALLOW responses for the dense attack distribution
of OPPORTUNISTIC traffic.
This interpretation is supported by the near-constant false-positive rate
of 1.0 throughout training: the DQN has effectively learned to never ALLOW,
which maximises reward on the training distribution but would fail on a more
balanced traffic mixture.

\subsection{Loss Dynamics}

The increasing loss trend through training (from $\approx$\,0.97 to
$\approx$\,2.4) is not indicative of divergence.
It reflects the expansion of the Q-value range as the policy learns to
differentiate high-reward from low-reward state-action pairs.
The Bellman target $y_i = r_i + \gamma \max_{a'} Q(s', a'; \theta^-)$
grows as the Q-function improves, causing the absolute Huber loss to
increase even as the relative TD error decreases.
The stable variance in the later training phase (no upward spikes after
episode 100) confirms that the stabilising mechanisms (target network,
gradient clipping) are effective.

\subsection{DQN Limitations in This Setting}

The DQN baseline reveals two limitations specific to the HoneyIQ setting.

\textbf{Class imbalance in training data.}
The OPPORTUNISTIC training episodes contain predominantly attack traffic.
The DQN is never exposed to episodes with significant benign traffic,
so it cannot learn to discriminate benign from malicious states.
The effective strategy of ``always respond with a non-ALLOW action'' achieves
high training reward but is not a robust policy.

\textbf{Transfer gap.}
Training exclusively on OPPORTUNISTIC intent means the Q-function is
calibrated to the specific attack distribution of OPPORTUNISTIC campaigns.
Evaluating against STEALTHY or TARGETED would require the DQN to generalise
to different threat-level time series; without intent-conditioned training
or multi-intent mixed training, this generalisation is likely to be
imperfect.
The SEDM avoids this limitation entirely through the analytical escalation
risk computation.

%% ── 5.3 Interpretability–Performance Trade-off ───────────────
\section{Interpretability and Performance Trade-off}
\label{sec:disc_tradeoff}

\subsection{The Case for SEDM}

The SEDM achieves $>$\,99\% detection across all four intent profiles with
zero training overhead, full transparency, and an auditable decision
procedure.
These properties are highly valuable for operational deployment:
%
\begin{itemize}
  \item \textbf{No training data requirement.}
        The SEDM requires only the Markov transition model, which encodes
        assumed attacker behaviour.
        It can be deployed immediately without a historical attack corpus.
  \item \textbf{Analytical escalation risk.}
        The escalation risk is computed from the transition matrix rather
        than estimated from data, making it exact under the model assumptions
        and independent of observation noise.
  \item \textbf{Full traceability.}
        Every action can be traced to a specific matrix cell and override
        rule.
        An analyst can reproduce any decision with pencil and paper.
  \item \textbf{Easy modification.}
        Security practitioners can update the matrix entries to reflect new
        threat intelligence (e.g., raising the response at WEAPONIZATION
        for a specific campaign) without retraining any model.
\end{itemize}

\subsection{The Case for DQN}

The DQN baseline offers adaptive flexibility that the SEDM cannot provide:
%
\begin{itemize}
  \item \textbf{Reward-optimised responses.}
        Given sufficient training data and a well-designed reward function,
        the DQN can discover optimal action strategies that are not
        immediately obvious to human designers.
  \item \textbf{Continuous adaptation.}
        Online DQN training (with a sliding replay buffer) allows the policy
        to adapt as the attack distribution evolves, provided new transitions
        are added to the buffer.
  \item \textbf{Nuanced feature sensitivity.}
        In a deployment with richer state representations (e.g., including
        classifier probability vectors, connection metadata), the DQN can
        exploit fine-grained patterns that a fixed matrix cannot capture.
\end{itemize}

\subsection{Practical Recommendation}

The results suggest a two-stage deployment strategy:
%
\begin{enumerate}
  \item \textbf{Initial deployment}: Use the SEDM as the primary policy.
        Its immediate operability, interpretability, and strong baseline
        performance ($>$\,99\% detection) make it the appropriate default.
  \item \textbf{Refinement}: Train a DQN agent alongside the SEDM using
        live honeypot interactions.
        If the DQN demonstrably outperforms the SEDM on detection rate or
        false positive rate under rigorous evaluation, and its decisions
        can be explained post-hoc via SHAP~\citep{lundberg2017unified} or
        similar techniques, it can be used to augment or eventually replace
        the SEDM for specific intent profiles.
\end{enumerate}

This strategy mirrors the human-AI collaboration model recommended by
\citet{rudin2019stop} for high-stakes decision-making: start with an
interpretable model, and only adopt a black-box model when there is
clear evidence of superior performance and adequate explanation mechanisms.

%% ── 5.4 Limitations ──────────────────────────────────────────
\section{Limitations}
\label{sec:disc_limitations}

\paragraph{Synthetic data only.}
Feature distributions are parametric approximations designed to reproduce
the qualitative signatures of UNSW-NB15 attack types.
They do not capture long-range temporal correlations, protocol interactions,
or the noise characteristics of real network traffic.
The simulation-to-reality gap means that results cannot be directly
extrapolated to live network deployments without validation on real
honeypot logs.

\paragraph{Defender action does not affect attacker trajectory.}
In the current formulation, the attacker's Markov chain evolves
independently of the defender's actions.
In practice, BLOCK or ALERT responses may deter, delay, or redirect an
attacker.
Modelling the attacker as a best-responding agent --- using multi-agent RL,
game-theoretic formulations, or a secondary attacker Markov chain that
conditions on the observed defender actions --- would produce more realistic
adversarial dynamics.

\paragraph{Classifier not integrated into the decision pipeline.}
The Random Forest classifier correctly identifies attack types from network
features, but its output is not used by the SEDM, which reads ground-truth
labels from the state vector.
Replacing ground-truth labels with classifier predictions would test whether
the SEDM is robust to classification noise --- an important prerequisite for
real deployment.
The classifier's misclassification rate (if non-trivial) could degrade the
escalation risk computation, which currently assumes that the attack type
is known.

\paragraph{No network topology.}
HoneyIQ models all interactions as single-session point contacts without
spatial structure.
Real network intrusions involve lateral movement, multi-hop attacks, and
topological constraints that are not captured.
Extending the simulation to a graph-based network topology would increase
realism at the cost of significantly higher state-space complexity.

\paragraph{Static Markov transition matrices.}
The attacker's transition matrices are fixed for each intent and do not
adapt during an episode.
Real adversaries update their strategies based on the defender's responses.
An adversary that observes consistent BLOCK actions on EXPLOITS might switch
to a different attack type or modify its escalation speed.
A dynamic Markov model whose transition matrices evolve in response to
defender feedback would capture this adaptive behaviour.

\paragraph{Episode termination by truncation.}
Episodes terminate only when \texttt{max\_steps} is reached; there is no
natural termination condition corresponding to attacker success or defeat.
In a real scenario, an attacker who achieves Actions on Objectives (data
exfiltration, system destruction) would terminate the campaign.
Modelling goal-conditioned episode termination would allow the evaluation
to capture the probability of attacker success, a more operationally
relevant metric than episode reward.

%% ── 5.5 Design Choices ───────────────────────────────────────
\section{Design Choices and Trade-offs}
\label{sec:disc_design}

\paragraph{LayerNorm over BatchNorm.}
LayerNorm~\citep{ba2016layer} normalises activations across features for
each sample independently, making it batch-size-agnostic.
For DQN training with mini-batches of 64 samples drawn from a diverse
replay buffer, BatchNorm statistics would be computed from a
heterogeneous mix of state types, potentially introducing instability.
LayerNorm avoids this dependency.
Empirically, no training instability attributable to normalisation was
observed during the 300-episode training runs.

\paragraph{Hard target network update.}
A hard copy of the policy network parameters every 150 steps was chosen
over a Polyak soft update ($\theta^- \leftarrow \tau\theta + (1-\tau)\theta^-$,
$\tau \approx 0.005$) for implementation simplicity.
Both approaches stabilise training; the hard update introduces discrete
target shifts every 150 steps but avoids the hyperparameter sensitivity
of choosing $\tau$.
At the episode length of 500 steps, 150-step updates correspond to roughly
three target updates per episode, providing stable targets over a substantial
fraction of each episode.

\paragraph{Composite threat-level weight assignment.}
The weight assignment in Equation~\ref{eq:threat} (attack severity 45\%,
kill chain stage 35\%, escalation rate 15\%, cumulative count 5\%) was
designed by domain reasoning rather than optimisation.
Attack severity receives the highest weight because it is the most direct
indicator of the attacker's current capability and impact.
Kill chain stage receives the second-highest weight because it captures the
campaign context.
The 5\% weight on cumulative attack count is a minor adjustment that
differentiates isolated probes from sustained campaigns without dominating
the threat signal.

An alternative approach would optimise these weights jointly with the reward
matrix using Bayesian optimisation or multi-objective optimisation, with
detection rate and false-positive rate as objectives.
This is left for future work.

\paragraph{SEDM band thresholds.}
The Low/Medium ($\rho = 0.35$) and Medium/High ($\rho = 0.65$) thresholds
were chosen to create roughly equal-width bands in the $[0, 1]$ risk space.
Intent-specific analysis (Figure~\ref{fig:esc_risk_per_intent}) confirms
that the thresholds correctly separate the intent profiles: STEALTHY
typically falls in the Low band, AGGRESSIVE in the High band, and
TARGETED/OPPORTUNISTIC in the Medium band for most stages.
Shifting the thresholds could alter the balance between false positives
(lower thresholds) and false negatives (higher thresholds) for specific
stages.

\paragraph{Synthetic classifier training data.}
Generating classifier training data from the parametric feature distributions
avoids the need for a real labelled dataset and enables instant re-training
with any desired class balance.
The risk is classifier over-fit to synthetic data that may not generalise
to real traffic.
In the HoneyIQ evaluation, this risk is moot because the evaluation features
are also synthetic; in a real deployment, the classifier would need to be
retrained on real network data.

%% ── 5.6 Future Work ──────────────────────────────────────────
\section{Future Work Directions}
\label{sec:future}

\subsection{Classifier Integration}
\label{sec:future_classifier}

The most immediate extension is to replace the ground-truth attack label
in the SEDM's decision pipeline with the Random Forest classifier's output.
This would introduce observation noise proportional to the classifier's
error rate and test whether the SEDM maintains $>$\,99\% detection under
realistic classification uncertainty.
A confidence-weighted band assignment --- where uncertain classifier
predictions widen the escalation risk estimate --- could provide a principled
way to handle ambiguous observations.

\subsection{Adaptive Attacker Modelling}

Modelling the attacker as a best-responding agent (e.g., using multi-agent
RL or a game-theoretic Stackelberg formulation) would introduce genuine
adversarial dynamics.
An attacker trained against the SEDM might learn to manipulate the
escalation risk computation by dwelling at low-risk stages longer than
the Markov chain predicts, or by generating traffic that bypasses the
override rules.
Evaluating SEDM robustness against such an adaptive adversary would provide
a stronger performance guarantee.

\subsection{Multi-Intent Training for DQN}

Exposing the DQN to a mixture of intent profiles during training --- sampling
a random intent at the start of each episode --- would broaden the state
distribution and reduce the tendency to over-specialise on OPPORTUNISTIC
traffic.
An intent-conditioned DQN, where the current intent is provided as input
(as it is in the 24-dimensional state vector), could be compared with the
SEDM on cross-intent generalisation.

\subsection{Real Network Feature Distributions}

Fitting the 15 feature distributions to the actual UNSW-NB15 records
(using maximum likelihood estimation or kernel density estimation per
attack category) would reduce the simulation-to-reality gap.
Alternatively, a conditional generative model (e.g., a conditional VAE
or normalising flow) trained on real packet data would enable richer,
correlated feature generation.

\subsection{Network Topology Extension}

Extending HoneyIQ to a multi-node network topology with explicit honeypot
placement decisions would introduce a spatial dimension to the problem.
The SEDM could be extended to consider both the kill chain stage and the
network position of the attacker, enabling topology-aware response policies.

\subsection{Distributional Reinforcement Learning}

Replacing the standard DQN with a distributional RL variant
(e.g., C51~\citep{wang2016dueling} or QR-DQN) would allow the agent to
maintain a distribution over returns rather than a point estimate.
The resulting uncertainty estimates could improve decisions at threat-band
boundaries, where the optimal action is least clear.

\subsection{OpenCanary Integration}

HoneyIQ includes initial scaffolding for integration with OpenCanary,
an open-source honeypot framework.
Replaying real OpenCanary session logs through the SEDM would provide an
out-of-distribution test of generalisation to genuine attack traffic,
bridging the simulation-to-reality gap without requiring live deployment.
